% =====================================================================
%  Interlude --- Shannon  (Scenes 13--16)
%  Pure section content. No preamble. LuaLaTeX-compatible.
% =====================================================================

% Interlude is austere: chrome accent goes to ink, both Act accents
% absent. The existing Shannon portrait divider below is the opener.
\AINITinterlude

% ---------------------------------------------------------------------
% Section divider --- opens the interlude
% ---------------------------------------------------------------------
\sectiondividerwithportrait%
  {shannon.jpg}%
  {``The fundamental problem of communication is that of\\ reproducing at one point either exactly or approximately\\ a message selected at another point.''}%
  {Claude E.\ Shannon, 1948}%
  {Interlude}

% ---------------------------------------------------------------------
% Scene 13 --- Shannon, 1948
% ---------------------------------------------------------------------
\begin{frame}{Shannon, 1948}
  \begin{columns}[T]
    \column{0.74\textwidth}
    In 1948, Claude Shannon published \AINITemph{A Mathematical Theory of
    Communication} --- a paper with two theorems at its core. The
    \AINITemph{source coding theorem} fixes a floor on how compactly any
    sequence of symbols can be described; the \AINITemph{channel coding
    theorem} fixes a ceiling on how reliably they can be transmitted
    through noise. For the next seven minutes we restate the first of
    these in the only vocabulary this lecture admits: urns.

    \column{0.22\textwidth}
    \vspace{0.2em}
    \includegraphics[width=\linewidth,keepaspectratio]{shannon.jpg}
  \end{columns}
\end{frame}

% ---------------------------------------------------------------------
% Scene 14 --- Entropy Mountain
%   Optional framing slide, then the viz handoff.
% ---------------------------------------------------------------------
\begin{frame}{How unpredictable is the next ball?}
  An urn with bias $p$ produces a stream of white and yellow draws.
  When $p = \tfrac{1}{2}$ every draw is a coin flip --- maximal
  surprise. When $p = 0.99$ almost every draw is white --- the
  sequence is nearly redundant. \AINITemph{Entropy} is the name we
  give to this gradient between certainty and surprise. Let us feel
  it before we write it down.
\end{frame}

\vizhandoff{entropy-mountain.html}

% ---------------------------------------------------------------------
% Scene 15 --- H(p) floor stated
% ---------------------------------------------------------------------
\equationslide%
  {The entropy of an urn}%
  {H(p) = -p\log_2 p - (1-p)\log_2 (1-p)}%
  {No compression scheme can describe an urn's draws in fewer than
  $H(p)$ bits per ball. Eckart--Young was the floor for compressing a
  matrix; entropy is the floor for compressing a sequence.}

% ---------------------------------------------------------------------
% Scene 16 --- Bridge into Act II
% ---------------------------------------------------------------------
% Quiet bridge slide. The "Act II --- Prediction" label is dropped
% because the Laboratory plate at the top of act2_prediction.tex now
% serves as the act opener; this slide stays purely as the synthesis
% sentence.
\sectiondivider%
  {\AINITemph{Recsys was compression. Bandits will be prediction.\\
  Same urns, different question.}}%
  {}%
  {}
